\documentclass[10pt,a4paper]{amsart}
\usepackage[margin=19mm]{geometry}
\usepackage{amsmath,amssymb,mathtools}
\usepackage[scr=rsfs,cal=euler]{mathalfa}
\usepackage{xfrac}
\usepackage[unicode,hidelinks,bookmarks=false]{hyperref}
\newcommand{\ie}{i.e.\ }
\newcommand{\cf}{cf.\ }
\newcommand{\resp}{respectively\ }
\newcommand{\Subsection}{\subsection}
\providecommand{\nobreakdash}{\nobreak\hbox{-}}
\setlength{\emergencystretch}{2em}
\multlinegap0pt
\usepackage{bm}
\usepackage{amssymb}
\usepackage{tikz}
\usepackage{enumerate}
\usepackage{cancel}
\usepackage{stackengine}
\newtheorem{lemma}{Lemma}
\newtheorem{remark}{Remark}
\newtheorem{theorem}{Theorem}
\newcommand{\FF}[1]{\underline{\overline{\rm{#1}}}}
\newcommand{\II}{{\rm  I\hspace{-.2mm}I}}
\renewcommand\geq{\geqslant}
\renewcommand\leq{\leqslant}
\title{{Higher Fundamental Forms and Warped Product Hypersurfaces }}
\author{Samuel Blitz${}^\flat$ \& Josef \v{S}ilhan${}^\sharp$}
\date{Source: arXiv:2410.06706v3 (CC BY 4.0)}
\begin{document}
\maketitle

\noindent\emph{Source and licence.} {Higher Fundamental Forms and Warped Product Hypersurfaces }. Version arXiv:2410.06706v3 (CC BY 4.0), distributed by arXiv under the Creative Commons Attribution 4.0 International licence, http://creativecommons.org/licenses/by/4.0/, as recorded in the arXiv licence metadata for this exact version and shown on the abstract page https://arxiv.org/abs/2410.06706v3. Full licence notice: \texttt{licenses/arxiv-2410.06706-CC-BY-4.0.txt}.\par\smallskip

\noindent\textit{Editorial scope.} Counted unit: the article's theorem vanishing-evens (source0.tex lines 622--624). The article's complete original proof of it is delivered with it (source0.tex lines 625--681, one \texttt{proof} environment of 57 lines). No mathematics is written, repaired or re-derived: the statement and the proof are the article's own lines, in the article's own notation, and no retained line is substituted or re-flowed.  Retained uncounted as the source-local context the proof needs: lines 408--412; lines 431--435; lines 448--452.  Nothing was omitted from any retained span, so the package carries no omission record.  No retained line contains a citation, so the wrapper carries no bibliography and no bibliography entry.  No retained line contains a figure environment, an image-inclusion command or drawing code, so no image is delivered and the package declares no asset dependency.  Editorial formatting only, in addition: an AMS \texttt{amsart} preamble that restates the article's own declarations and macros used by the retained lines, namely the environments \texttt{lemma}, \texttt{remark}, \texttt{theorem} on separate counters, together with the standard packages those lines need.  The retained environments print in this document as Lemma 1, Remark 1, Theorem 1 in order of appearance, whereas the article prints the same items with the numbering of its own full document.  The complete native source of the article as distributed in the arXiv e-print for this exact version is bundled unchanged as \texttt{sources/166538/source0.tex}.
\begin{lemma} \label{normal-derivs}
Let $\iota : \Sigma \hookrightarrow (M,g)$ be a smooth embedded hypersurface with the vector field $n$ dual to 
the unit conormal 1-form. Then, for every $k \geq 2$, we have that
$$\FF{k} = \iota^* \nabla_n^{k-2} \nabla n + \text{products of lower order fundamental forms.}$$
\end{lemma}

\begin{lemma} \label{lie-derivs}
Let $\iota : \Sigma \hookrightarrow (M,g)$ be a smooth embedded hypersurface with the vector field $n$ dual to 
the unit conormal 1-form. Then, for every $k \geq 2$, we have that
$$\FF{k} = \tfrac{1}{2} \iota^* \mathcal{L}_n^{k-1} g+ \text{products of lower order fundamental forms.}$$
\end{lemma}

\begin{remark}\label{weights}
Note that in all of these constructions involving equality modulo lower order terms, we may ascribe weights to better control which fundamental forms appear in each identity. Indeed, assigning a weight of $3-k$ to $\FF{k}$ and a weight of $-2$ for every contraction, it is clear that the weights of each summand in the above identities are always equal.

Additionally, observe that when the lower fundamental forms vanish, the original definition of the fundamental forms is equivalent to the constructions provided in Lemmas~\ref{normal-derivs} and~\ref{lie-derivs}.
\end{remark}

\begin{theorem} \label{vanishing-evens}
Let $\Sigma \hookrightarrow (M,g)$ be a base-like hypersurface embedding with warping function $f$. Then, $\FF{k} = 0$ for all $k = 2n$ for positive integers $n$.
\end{theorem}

\begin{proof}
As $(M,g)$ is a warped product, in a coordinate system adapted to the warped product $(t, x^i)$, the metric is expressible as
$$g = f(x^1, \ldots, x^{d-1})^2 dt^2 + \bar{g}_{ij}(x^1, \ldots, x^{d-1}) dx^i dx^j\,.$$
%Now, we may explicitly compute the second fundamental form in this metric, noting that $n = \sqrt{f} dt$:
%$$\iota^* \nabla_a n_b = \iota^* [(dt)_b \partial_a \sqrt{f}] - \sqrt{f} \iota^* \Gamma^t_{ab}\,.$$
%An explicit computation of the required Christoffel symbols show that they vanish, so we have that $\II = 0$.

Now we wish to change coordinates so that $g$ is asymptotically in normal form, i.e. so that
$$g = F^{(n+2)}(s,y^j) ds^2 + G_i^{(n+1)}(s,y^j) ds dy^i + \gamma_{ij}(s,y^1,\ldots,y^{d-1}) dy^i dy^j\,,$$
where $n \geq 0$ is even and $F^{(n+2)}(s,y^j) = 1 + \mathcal{O}(s^{n+2})$ is an even function of $s$. Similarly, we will require that $G_i^{(n+1)} = \mathcal{O}(s^{n+1})$ is an odd function of $s$.

To find such a coordinate transformation, we solve order-by-order inductively. The base case is fairly easy, with $n=0$: define $t = \frac{s}{f(y^i)}$ and $x^i = y^i$. Then, we have that
$$g = ds^2 - \frac{2s \partial_{y^i} f}{f}  dy^i ds + \gamma_{ij} dy^i dy^j \,,$$
for some $\gamma_{ij}$. Clearly, here $F^{(2)}(s,y^j) = 1$ is even and $G_i^{(1)}(s,y^j) = -\frac{2s \partial_{y^i} f}{f}$ is odd in $s$.
We now use induction to proceed. 

Suppose that $n \geq 0$ is even, that there exists $T_{n+1}(s,y^j) = \frac{s}{f(y^j)} + \cdots + Q_{n+1}(y^j) s^{n+1}$ an odd polynomial in $s$, and $X^i_{n}(s,y^j) = y^i + \cdots + P^i_{n}(y^j) s^{n}$ an even function of $s$ such that
$$g = F^{(n+2)}(s,y^j) ds^2 + G_i^{(n+1)}(s,y^j) ds dy^i + \gamma_{ij} dy^i dy^j.$$
Here we use coordinates $(s,y^j)$ related to the original coordinates $(t,x^i)$ by
$t = T_{n+1}(s,y^j)$ and $x^i = X^i_{n}(s,y^j)$, where $F^{(n+2)}$ is an even function of $s$ (with first coefficient 1)
and $G_i^{(n+1)}$ is an odd function of $s$. When $n=0$, this is the base case.

We now increment $n \rightarrow n+2$. Let $t = T_{n+1}(s,y^j) + s^{n+3} A_{n+3}(y^j)$ and $x^i = X^i_{n}(s,y^j) + s^{n+2} B^i_{n+2}(y^j)$, with $A_{n+3}$ and $B^i_{n+2}$ arbitrary functions of $\vec{y}$. We would like to show that there exist unique $A_{n+3}$ and $B^i_{n+2}$ such that 
$g_{ss} = F^{(n+4)} = 1 + \mathcal{O}(s^{n+4})$ and $g_{si} = G^{(n+3)}_i = \mathcal{O}(s^{n+3})$. 
To do so, we may compute components of the metric:
\begin{align*}
g_{si} =& G^{(n+1)}_i +  (n+2)  s^{n+1} \bar{g}_{jk} (\partial_{y^i} X_n^j) B^k_{n+2} \\
&+ s^{n+2} \left((n+3) g_{tt} A_{n+3} \partial_{y^i} T_{n+1} + \bar{g}_{jk} (\partial_s X^j_n)\partial_{y^i} B^k_{n+2} \right) \\
&+ s^{n+3} g_{tt} (\partial_s T_{n+1}) \partial_{y^i} A_{n+3} \\
&+ (n+2) s^{2n+3} \bar{g}_{jk} B^j_{n+2} \partial_{y^i} B^k_{n+2} + (n+3) s^{2n+5} g_{tt} A_{n+3}  \partial_{y^i} A_{n+3} \,.
\end{align*}
Because $\bar{g}_{jk}$ is invertible, we may fix $B^i_{n+2}$ so that the second term cancels with the leading term of $G^{(n+1)}_i$. Because $G^{(n+1)}_i$ is odd and $X^j_n$ is even (satisfying $\partial_{y^i} X^j_n = \delta_i^j + \mathcal{O}(s^2)$), the remaining higher-order terms are shunted to order $n+3$ and remain odd. Then, note that because both $T_{n+1}$ and $\partial_s X^j_n$ are odd in $s$, we see that the second line is actually odd and is of order $s^{n+3}$. Indeed, it follows that $g_{si}$ is odd and of order $s^{n+3}$. Now having fixed $B^i_{n+2}$, we may consider $g_{ss}$:
\begin{align*}
g_{ss} =& 1 + 2(n+2) s^{n+1} \bar{g}_{jk} B^j_{n+2} \partial_s X^k_n \\
&+ (F^{(n+2)} - 1)  + 2(n+3) s^{n+2} g_{tt} A_{n+3} \partial_s T_{n+1} \\
&+ (n+2)^2 s^{2n+2} \bar{g}_{jk} B^j_{n+2} B^k_{n+2} + (n+3)^2 s^{2n+4} g_{tt} A_{n+3}^2 \,.
\end{align*}
Now because $B^i_{n+2}$ is fixed and $\partial_s X^k_n$ is odd, the second term on the first line is 
even and has the factor $s^{n+2}$. The same holds for $F^{(n+2)} - 1$ and also for the second
term on the middle lide of the previous display since $T_{n+1}$ is odd. Since moreover $\partial_s T_{n+1}|_{s = 0} \neq 0$, 
we may choose $A_{n+3}$ uniquely so that the sum of first two lines of the previous display is of the form
$1 + \mathcal{O}(s^{n+4})$. 
But again, $\partial_s T_{n+1}$ is even as is $F^{(n+2)}$, so all higher-order terms remain even. Thus, the induction holds. 
Note $\gamma_{jk}$ changes in an induction step but since we did not impose any particular condition
on these functions, we keep the same notation $\gamma_{jk}$ for all induction steps.


Given such a coordinate transformation, we can now determine the parity of $\gamma_{ij}$. However, this is trivially the case: $\frac{dt}{dy^i}$ is odd, so $\frac{dt}{dy^i} \frac{dt}{dy^j}$ is even. Similarly, $\frac{dx^a}{dy^i} \frac{dx^b}{dy^j}$ is even. Therefore, we find that $\gamma_{ij}$ is an even function of $s$. Taking stock of this coordinate transform, we have that
$$g = ds^2 + (\gamma^{(0)}_{ij} + s^2 \gamma^{(2)}_{ij} + \cdots) dy^i dy^j + \mathcal{O}(s^m)$$
for any choice of $m$. Therefore, we have that in these coordinates, $n^a =  \partial_s + \mathcal{O}(s^m)$, and so $\mathcal{L}_{n} T = \partial_s T + \mathcal{O}(s^{m-1})$ for any tensor $T$. So, for any $k \leq m$, we have that $\iota^* \mathcal{L}_n^k g = \iota^* \partial_s^k g$. Notably, for any $k$ odd, we have that we may always find a coordinate transform so that
$$\iota^* \partial_s^k g = 0\,.$$

Now recall from Lemma~\ref{lie-derivs} that $\iota^* \mathcal{L}_n^k g$ is proportional to the $(k+1)$th fundamental form, modulo lower order fundamental forms. So it follows directly from  Lemma~\ref{lie-derivs} that when $k$ is odd,
$$ \FF{k+1} = \text{products of lower order fundamental forms.}$$
Here, as per Remark~\ref{weights}, each summand on the right hand side of the above display has odd weight $2-k$. But odd fundamental forms always have even weight, as does $\bar{g}^{-1}$. It thus follows that no summand on the right hand side of the above display is composed entirely of odd fundamental forms. Hence, it follows that if $\FF{m}$ vanishes for each $2 \leq m \leq k-1$ even, then it also follows that $\FF{k+1}$ vanishes. This completes the induction and hence the proof.

\end{proof}

\end{document}
